\documentclass[11pt,a4paper]{article}
\usepackage[utf8]{inputenc}
\usepackage[italian]{babel}
\usepackage{geometry}
\usepackage{graphicx}
\usepackage{hyperref}
\usepackage{listings}
\usepackage{xcolor}
\usepackage{amsmath}
\usepackage{enumitem}
\usepackage{booktabs}
\usepackage{longtable}
\usepackage{array}
\usepackage{tcolorbox}

\geometry{margin=2.5cm}

% Configurazione codice
\lstset{
    language=JavaScript,
    basicstyle=\ttfamily\small,
    keywordstyle=\color{blue}\bfseries,
    stringstyle=\color{red},
    commentstyle=\color{green!60!black},
    numberstyle=\tiny\color{gray},
    numbers=left,
    numbersep=5pt,
    frame=single,
    breaklines=true,
    showstringspaces=false,
    tabsize=2,
    escapeinside={(*@}{@*)}
}

% Box per evidenziare sezioni importanti
\newtcolorbox{importantbox}[1]{
    colback=red!5!white,
    colframe=red!75!black,
    title=#1,
    fonttitle=\bfseries
}

\newtcolorbox{warningbox}[1]{
    colback=orange!5!white,
    colframe=orange!75!black,
    title=#1,
    fonttitle=\bfseries
}

\newtcolorbox{tipbox}[1]{
    colback=blue!5!white,
    colframe=blue!75!black,
    title=#1,
    fonttitle=\bfseries
}

\title{\textbf{Analisi Componente: BaseService}\\
\large Documentazione Tecnica Esaustiva}
\author{Documentazione Sistema}
\date{\today}

\begin{document}

\maketitle
\tableofcontents
\newpage

\section{Responsabilità}

\subsection{Cosa fa esattamente questo componente?}

\textbf{BaseService} è una classe astratta che implementa operazioni CRUD (Create, Read, Update, Delete) standardizzate per tutti i modelli del sistema, con \textbf{isolamento multi-tenant esplicito e sicuro}.

\textbf{Funzioni principali}:
\begin{itemize}
    \item \textbf{Astrazione dell'accesso ai dati}: Il controller non conosce i dettagli di MongoDB
    \item \textbf{Multi-tenancy automatico}: Ogni query include automaticamente il filtro \texttt{\{user: userId\}}
    \item \textbf{Prevenzione IDOR (Insecure Direct Object Reference)}: Verifica ownership esplicita su ogni operazione
    \item \textbf{Prevenzione Mass Assignment}: Blocca modifica di campi sensibili (\texttt{user}, \texttt{\_id})
    \item \textbf{Lifecycle Hooks}: Permette alle sottoclassi di estendere il comportamento senza modificare il codice base
\end{itemize}

\subsection{Perché esiste?}

\begin{importantbox}{Problema Risolto}
Senza BaseService, ogni service dovrebbe:
\begin{itemize}
    \item Ripetere la logica di filtraggio per \texttt{user} in ogni metodo
    \item Implementare manualmente la verifica di ownership
    \item Gestire errori comuni (documento non trovato, validazione, etc.)
    \item Scrivere codice duplicato per CRUD standard
\end{itemize}
\end{importantbox}

\textbf{Vantaggi dell'approccio}:
\begin{enumerate}
    \item \textbf{DRY (Don't Repeat Yourself)}: Logica comune centralizzata
    \item \textbf{Sicurezza by Default}: Impossibile dimenticare il filtro \texttt{user}
    \item \textbf{Consistenza}: Tutti i service si comportano allo stesso modo
    \item \textbf{Manutenibilità}: Modifiche alla logica comune in un solo posto
    \item \textbf{Testabilità}: Facile testare la logica base una volta
\end{enumerate}

\section{Struttura del Codice}

\subsection{Architettura}

BaseService implementa due pattern principali:

\begin{enumerate}
    \item \textbf{Template Method Pattern}: I metodi base definiscono lo "scheletro" dell'algoritmo, mentre le sottoclassi possono sovrascrivere hook specifici (\texttt{beforeCreate}, \texttt{afterUpdate}, etc.)
    \item \textbf{Repository Pattern (semplificato)}: Astrae l'accesso ai dati. Il controller non sa se i dati vengono da MongoDB, PostgreSQL o un'API esterna.
\end{enumerate}

\subsection{Metodi Principali}

\subsubsection{Helper Methods (Private)}

\begin{lstlisting}
// Estrae userId dal tenantScope con validazione fail-fast
_getUserId(tenantScope) {
    if (!tenantScope || !tenantScope.tenantId) {
        throw new Error('SECURITY VIOLATION: tenantScope.tenantId mancante!');
    }
    return tenantScope.tenantId;
}

// Crea filtro base con user incluso
_buildFilter(tenantScope, additionalFilters = {}) {
    const userId = this._getUserId(tenantScope);
    return { user: userId, ...additionalFilters };
}
\end{lstlisting}

\textbf{Logica}: 
\begin{itemize}
    \item \texttt{\_getUserId()}: Validazione fail-fast. Se \texttt{tenantScope} è null o manca \texttt{tenantId}, lancia errore immediatamente (previene data leak silenziosi)
    \item \texttt{\_buildFilter()}: Combina sempre \texttt{\{user: userId\}} con filtri aggiuntivi. Garantisce che ogni query sia isolata per tenant
\end{itemize}

\subsubsection{Read Operations}

\begin{lstlisting}
// Trova tutti i documenti del tenant
async find(tenantScope, filters = {}, options = {}) {
    const secureFilter = this._buildFilter(tenantScope, filters);
    let query = this.Model.find(secureFilter);
    
    // Applica sort, limit, skip, select, populate
    if (sort) query = query.sort(sort);
    if (limit) query = query.limit(limit);
    // ...
    
    return query.exec();
}

// Trova documento per ID (con verifica tenant)
async findById(tenantScope, id, options = {}) {
    // ⚠️ USA findOne con { _id, user } invece di findById
    // Questo previene IDOR anche se l'attaccante conosce l'ID
    const secureFilter = this._buildFilter(tenantScope, { _id: id });
    const doc = await this.Model.findOne(secureFilter).exec();
    
    if (!doc && throwIfNotFound) {
        throw AppError.notFound(this.options.entityName);
    }
    return doc;
}
\end{lstlisting}

\textbf{Logica}:
\begin{itemize}
    \item \texttt{find()}: Aggiunge automaticamente \texttt{\{user: userId\}} alla query. Supporta sort, limit, skip, select, populate
    \item \texttt{findById()}: \textbf{Non usa} \texttt{Model.findById()} perché non verifica ownership. Usa \texttt{findOne(\{_id, user\})} per prevenire IDOR
    \item \texttt{findOne()}: Trova un documento con filtri custom, sempre con filtro \texttt{user}
    \item \texttt{count()}: Conta documenti con filtro tenant
    \item \texttt{exists()}: Verifica esistenza documento
\end{itemize}

\subsubsection{Write Operations}

\begin{lstlisting}
// Crea nuovo documento
async create(tenantScope, data) {
    const userId = this._getUserId(tenantScope);
    
    // SICUREZZA: Rimuovi 'user' e '_id' dai dati in input
    // Prevenzione mass assignment attacks
    const { user, _id, ...safeData } = data;
    
    // Hook pre-create (override nelle sottoclassi)
    const processedData = await this.beforeCreate(tenantScope, safeData);
    
    // 🔒 SICUREZZA ESPLICITA: Iniettiamo userId direttamente
    const doc = await this.Model.create({
        ...processedData,
        user: userId, // ← ESPLICITO, non delegato a plugin
    });
    
    // Hook post-create
    await this.afterCreate(tenantScope, doc);
    
    return doc;
}

// Aggiorna documento per ID
async update(tenantScope, id, data, options = {}) {
    // Blocca modifica di campi sensibili
    const { user, _id, ...safeData } = data;
    
    // Hook pre-update
    const processedData = await this.beforeUpdate(tenantScope, id, safeData);
    
    // Filtro { _id, user } per prevenire IDOR
    const secureFilter = this._buildFilter(tenantScope, { _id: id });
    
    const doc = await this.Model.findOneAndUpdate(
        secureFilter,
        { $set: processedData },
        { new: true, runValidators: true, ...options }
    );
    
    if (!doc) {
        throw AppError.notFound(this.options.entityName);
    }
    
    await this.afterUpdate(tenantScope, doc);
    return doc;
}

// Elimina documento per ID
async delete(tenantScope, id) {
    // Hook pre-delete (es. verifica dipendenze)
    await this.beforeDelete(tenantScope, id);
    
    const secureFilter = this._buildFilter(tenantScope, { _id: id });
    const doc = await this.Model.findOneAndDelete(secureFilter);
    
    if (!doc) {
        throw AppError.notFound(this.options.entityName);
    }
    
    // Hook post-delete (es. cleanup)
    await this.afterDelete(tenantScope, doc);
    return doc;
}
\end{lstlisting}

\textbf{Logica}:
\begin{itemize}
    \item \texttt{create()}: Inietta \texttt{user: userId} esplicitamente. Rimuove \texttt{user} e \texttt{\_id} da \texttt{data} per prevenire mass assignment
    \item \texttt{update()}: Usa \texttt{findOneAndUpdate} con filtro \texttt{\{_id, user\}} per prevenire IDOR. Blocca modifica di \texttt{user} e \texttt{\_id}
    \item \texttt{delete()}: Verifica ownership prima di eliminare. Supporta hook pre/post per validazioni e cleanup
    \item \texttt{deleteMany()}: Richiede filtri espliciti (non può essere vuoto) per prevenire eliminazioni accidentali
\end{itemize}

\subsubsection{Utility Methods}

\begin{lstlisting}
// Ricerca testuale sui campi configurati
async search(tenantScope, searchTerm, additionalFilters = {}) {
    if (!searchTerm || !this.options.searchFields.length) {
        return this.find(tenantScope, additionalFilters);
    }
    
    const regex = new RegExp(searchTerm, 'i');
    const searchConditions = this.options.searchFields.map(field => ({
        [field]: regex,
    }));
    
    return this.find(tenantScope, {
        ...additionalFilters,
        $or: searchConditions,
    });
}

// Paginazione con metadata
async paginate(tenantScope, filters = {}, { page = 1, limit = 20 } = {}) {
    const skip = (page - 1) * limit;
    
    const [data, total] = await Promise.all([
        this.find(tenantScope, filters, { skip, limit }),
        this.count(tenantScope, filters),
    ]);
    
    return {
        data,
        meta: {
            page, limit, total,
            totalPages: Math.ceil(total / limit),
            hasMore: page * limit < total,
        },
    };
}
\end{lstlisting}

\textbf{Logica}:
\begin{itemize}
    \item \texttt{search()}: Crea regex case-insensitive e cerca su \texttt{searchFields} configurati. Usa \texttt{\$or} per cercare su più campi
    \item \texttt{paginate()}: Esegue \texttt{find} e \texttt{count} in parallelo con \texttt{Promise.all}. Ritorna dati + metadata (page, total, hasMore)
\end{itemize}

\subsubsection{Lifecycle Hooks}

\begin{lstlisting}
// Hook eseguiti PRIMA delle operazioni
async beforeCreate(tenantScope, data) { return data; }
async beforeUpdate(tenantScope, id, data) { return data; }
async beforeDelete(tenantScope, id) { /* override */ }

// Hook eseguiti DOPO le operazioni
async afterCreate(tenantScope, doc) { /* override */ }
async afterUpdate(tenantScope, doc) { /* override */ }
async afterDelete(tenantScope, doc) { /* override */ }
\end{lstlisting}

\textbf{Logica}: Le sottoclassi possono sovrascrivere questi hook per:
\begin{itemize}
    \item \textbf{beforeCreate/Update}: Validazioni custom, trasformazioni dati, calcoli automatici
    \item \textbf{afterCreate/Update}: Side effects (notifiche, logging, aggiornamento relazioni)
    \item \textbf{beforeDelete}: Verifica dipendenze (es. non eliminare progetto con worklog)
    \item \textbf{afterDelete}: Cleanup (es. eliminare file associati)
\end{itemize}

\section{Dipendenze}

\subsection{Dipendenze Dirette}

\begin{itemize}
    \item \textbf{AppError} (\texttt{../utils/AppError}): Classe custom per errori standardizzati
    \item \textbf{Model Mongoose}: Passato al constructor. BaseService è agnostico rispetto al modello specifico
\end{itemize}

\subsection{Dipendenze Indirette}

\begin{itemize}
    \item \textbf{Mongoose}: ODM per MongoDB (usato tramite Model)
    \item \textbf{tenantContext Middleware}: Deve essere applicato alle routes per iniettare \texttt{req.tenantScope}
\end{itemize}

\subsection{Utilizzo nelle Sottoclassi}

BaseService è esteso da:
\begin{itemize}
    \item \texttt{GoalService}: Gestione obiettivi
    \item \texttt{WorkLogService}: Gestione worklog
    \item \texttt{ProjectService}: Gestione progetti
    \item \texttt{CheckInService}: Gestione check-in obiettivi
    \item \texttt{WorkTodoService}: Gestione TODO workspace
    \item \texttt{StudyService}: Gestione deck flashcard
    \item \texttt{WorkProjectService}: Gestione progetti workspace
    \item \texttt{WorkEntryService}: Gestione work entries
\end{itemize}

\section{Analisi Critica}

\subsection{Punti di Forza}

\begin{enumerate}
    \item \textbf{Sicurezza by Default}:
    \begin{itemize}
        \item Ogni query include \texttt{\{user: userId\}} esplicitamente
        \item Fail-fast su \texttt{tenantScope} mancante
        \item Prevenzione IDOR con \texttt{findOne(\{_id, user\})} invece di \texttt{findById()}
        \item Prevenzione mass assignment (rimozione \texttt{user}, \texttt{\_id} da input)
    \end{itemize}
    
    \item \textbf{Template Method Pattern Ben Implementato}:
    \begin{itemize}
        \item Lifecycle hooks permettono estensione senza modificare codice base
        \item Metodi base definiscono algoritmo, sottoclassi personalizzano comportamento
    \end{itemize}
    
    \item \textbf{DRY (Don't Repeat Yourself)}:
    \begin{itemize}
        \item Logica comune centralizzata
        \item Riduce codice duplicato del 70-80\% nei service
    \end{itemize}
    
    \item \textbf{Consistenza}:
    \begin{itemize}
        \item Tutti i service si comportano allo stesso modo
        \item Facilita onboarding nuovi sviluppatori
    \end{itemize}
    
    \item \textbf{Testabilità}:
    \begin{itemize}
        \item Facile testare logica base una volta
        \item Sottoclassi testano solo logica custom
    \end{itemize}
\end{enumerate}

\subsection{Possibili Bug/Rischi}

\begin{warningbox}{Edge Cases Non Gestiti}
\end{warningbox}

\begin{enumerate}
    \item \textbf{tenantScope null/undefined}:
    \begin{itemize}
        \item \textbf{Stato attuale}: \texttt{\_getUserId()} lancia errore (fail-fast)
        \item \textbf{Rischio}: Errore generico, non AppError standardizzato
        \item \textbf{Miglioramento}: Lanciare \texttt{AppError.unauthorized()} invece di \texttt{Error} generico
    \end{itemize}
    
    \item \textbf{Concorrenza su update()}:
    \begin{itemize}
        \item \textbf{Stato attuale}: Usa \texttt{findOneAndUpdate} che è atomico
        \item \textbf{Rischio}: Se due richieste aggiornano lo stesso documento simultaneamente, l'ultima vince (last-write-wins)
        \item \textbf{Miglioramento}: Considerare optimistic locking con \texttt{versionKey} di Mongoose
    \end{itemize}
    
    \item \textbf{deleteMany() senza filtri}:
    \begin{itemize}
        \item \textbf{Stato attuale}: Richiede filtri espliciti (già gestito)
        \item \textbf{Rischio}: Nessuno (già protetto)
    \end{itemize}
    
    \item \textbf{Populate su campi non esistenti}:
    \begin{itemize}
        \item \textbf{Stato attuale}: Mongoose ignora silenziosamente populate su campi non definiti nello schema
        \item \textbf{Rischio}: Basso, ma potrebbe confondere durante il debug
        \item \textbf{Miglioramento}: Validare \texttt{populateFields} nel constructor
    \end{itemize}
    
    \item \textbf{search() con regex injection}:
    \begin{itemize}
        \item \textbf{Stato attuale}: \texttt{searchTerm} viene usato direttamente in regex
        \item \textbf{Rischio}: Se \texttt{searchTerm} contiene caratteri speciali regex (\texttt{.}, \texttt{*}, \texttt{+}, etc.), potrebbe causare comportamenti inattesi
        \item \textbf{Miglioramento}: Escapare caratteri speciali: \texttt{searchTerm.replace(/[.*+?^${}()|[\]\\]/g, '\\\$&')}
    \end{itemize}
    
    \item \textbf{Limit/Skip senza validazione}:
    \begin{itemize}
        \item \textbf{Stato attuale}: Accetta qualsiasi valore numerico
        \item \textbf{Rischio}: \texttt{limit: -1} o \texttt{limit: 999999} potrebbero causare problemi
        \item \textbf{Miglioramento}: Validare \texttt{limit} (min: 1, max: 100) e \texttt{skip} (min: 0)
    \end{itemize}
\end{enumerate}

\subsection{Performance}

\begin{enumerate}
    \item \textbf{Query con populate multipli}:
    \begin{itemize}
        \item \textbf{Problema}: Se \texttt{populateFields} contiene molti campi, ogni query fa multiple join
        \item \textbf{Impatto}: Query lente su collezioni grandi
        \item \textbf{Miglioramento}: Considerare \texttt{select} per limitare campi popolati, o usare aggregation pipeline
    \end{itemize}
    
    \item \textbf{paginate() con count su collezioni grandi}:
    \begin{itemize}
        \item \textbf{Problema}: \texttt{countDocuments()} su collezioni con milioni di documenti è lento
        \item \textbf{Impatto}: Paginazione lenta
        \item \textbf{Miglioramento}: Considerare stima approssimata con \texttt{estimatedDocumentCount()} o caching del totale
    \end{itemize}
    
    \item \textbf{search() con \$or su molti campi}:
    \begin{itemize}
        \item \textbf{Problema}: Se \texttt{searchFields} contiene 10+ campi, \texttt{\$or} diventa complesso
        \item \textbf{Impatto}: Query lente, non può usare indici efficientemente
        \item \textbf{Miglioramento}: Limitare \texttt{searchFields} a 3-5 campi più rilevanti, o usare MongoDB text search
    \end{itemize}
    
    \item \textbf{Mancanza di indici}:
    \begin{itemize}
        \item \textbf{Problema}: BaseService non crea indici automaticamente
        \item \textbf{Impatto}: Query lente su collezioni grandi
        \item \textbf{Miglioramento}: Documentare che i modelli devono avere indice su \texttt{\{user: 1\}} e \texttt{\{user: 1, \_id: 1\}}
    \end{itemize}
\end{enumerate}

\subsection{Miglioramenti Codice (Clean Code)}

\begin{tipbox}{Suggerimenti per Migliorare}
\end{tipbox}

\begin{enumerate}
    \item \textbf{Validazione Input}:
    \begin{lstlisting}
// PRIMA
async findById(tenantScope, id, options = {}) {
    const secureFilter = this._buildFilter(tenantScope, { _id: id });
    // ...
}

// DOPO
async findById(tenantScope, id, options = {}) {
    if (!id || !mongoose.Types.ObjectId.isValid(id)) {
        throw AppError.badRequest('ID non valido');
    }
    const secureFilter = this._buildFilter(tenantScope, { _id: id });
    // ...
}
    \end{lstlisting}
    
    \item \textbf{Error Handling Standardizzato}:
    \begin{lstlisting}
// PRIMA
_getUserId(tenantScope) {
    if (!tenantScope || !tenantScope.tenantId) {
        throw new Error('SECURITY VIOLATION...');
    }
    return tenantScope.tenantId;
}

// DOPO
_getUserId(tenantScope) {
    if (!tenantScope || !tenantScope.tenantId) {
        throw AppError.unauthorized(
            'Sessione non valida. Effettua nuovamente il login.',
            { code: 'INVALID_TENANT_SCOPE' }
        );
    }
    return tenantScope.tenantId;
}
    \end{lstlisting}
    
    \item \textbf{Escapare Regex in search()}:
    \begin{lstlisting}
// PRIMA
const regex = new RegExp(searchTerm, 'i');

// DOPO
const escapedTerm = searchTerm.replace(/[.*+?^${}()|[\]\\]/g, '\\$&');
const regex = new RegExp(escapedTerm, 'i');
    \end{lstlisting}
    
    \item \textbf{Validazione Limit/Skip}:
    \begin{lstlisting}
// PRIMA
if (limit) query = query.limit(limit);

// DOPO
if (limit) {
    const validLimit = Math.max(1, Math.min(100, parseInt(limit, 10) || 20));
    query = query.limit(validLimit);
}
    \end{lstlisting}
    
    \item \textbf{TypeScript per Type Safety}:
    \begin{itemize}
        \item Convertire BaseService in TypeScript per type checking
        \item Definire interfacce per \texttt{tenantScope}, \texttt{options}, etc.
        \item Migliora autocomplete e previene errori a compile-time
    \end{itemize}
    
    \item \textbf{Logging per Debug}:
    \begin{lstlisting}
// Aggiungere logging opzionale per debug
async find(tenantScope, filters = {}, options = {}) {
    const secureFilter = this._buildFilter(tenantScope, filters);
    
    if (process.env.DEBUG_SERVICES === 'true') {
        console.log(`[BaseService.find] Model: ${this.Model.modelName}, Filter:`, secureFilter);
    }
    
    // ...
}
    \end{lstlisting}
\end{enumerate}

\section{UI/UX Impact (Indiretto per Backend)}

\subsection{Velocità di Risposta}

\begin{itemize}
    \item \textbf{Query Ottimizzate}: BaseService supporta \texttt{limit}, \texttt{skip}, \texttt{select} per ridurre dati trasferiti
    \item \textbf{Paginazione}: \texttt{paginate()} permette UI con infinite scroll o paginazione classica
    \item \textbf{Populate Selettivo}: Evita over-fetching di dati non necessari
\end{itemize}

\subsection{Sicurezza dei Dati}

\begin{importantbox}{Impatto Sicurezza}
BaseService garantisce che:
\begin{itemize}
    \item \textbf{Isolamento Completo}: Utente A non può mai vedere/modificare dati di Utente B
    \item \textbf{Prevenzione IDOR}: Anche conoscendo l'ID di un documento, non si può accedere se non si è il proprietario
    \item \textbf{Prevenzione Mass Assignment}: Campi sensibili (\texttt{user}, \texttt{\_id}) non possono essere modificati dall'input utente
\end{itemize}
\end{importantbox}

\subsection{Messaggi di Errore}

\begin{itemize}
    \item \textbf{Standardizzati}: Tutti i service usano \texttt{AppError} per messaggi consistenti
    \item \textbf{User-Friendly}: \texttt{AppError.notFound()} genera messaggi leggibili ("Obiettivo non trovato" invece di "Document not found")
    \item \textbf{Localizzazione}: \texttt{entityName} permette messaggi personalizzati per tipo di entità
\end{itemize}

\subsection{Esperienza Sviluppatore}

\begin{itemize}
    \item \textbf{API Consistente}: Tutti i service hanno gli stessi metodi (\texttt{find}, \texttt{create}, \texttt{update}, \texttt{delete})
    \item \textbf{Documentazione Implicita}: Se conosci BaseService, conosci tutti i service
    \item \textbf{Onboarding Veloce}: Nuovi sviluppatori capiscono subito come funzionano i service
\end{itemize}

\section{Esempi di Utilizzo}

\subsection{Esempio 1: GoalService (Sottoclasse Semplice)}

\begin{lstlisting}
class GoalService extends BaseService {
    constructor() {
        super(Goal, {
            searchFields: ['title', 'description'],
            defaultSort: { deadline: 1 },
            entityName: 'Obiettivo',
        });
    }
    
    // Metodo custom
    async findActive(tenantScope) {
        return this.find(tenantScope, { status: 'active' });
    }
}

// Utilizzo nel controller
const goals = await goalService.findActive(req.tenantScope);
// Query eseguita: Goal.find({ user: userId, status: 'active' })
\end{lstlisting}

\subsection{Esempio 2: ProjectService (Sottoclasse con Hook)}

\begin{lstlisting}
class ProjectService extends BaseService {
    constructor() {
        super(Project, {
            searchFields: ['name', 'code'],
            defaultSort: { name: 1 },
            entityName: 'Progetto',
        });
    }
    
    // Hook pre-create: calcola health del progetto
    async beforeCreate(tenantScope, data) {
        if (data.estimatedHours) {
            data.health = this._calculateHealth(data);
        }
        return data;
    }
    
    // Hook post-delete: verifica dipendenze
    async beforeDelete(tenantScope, id) {
        const workLogCount = await WorkLog.countDocuments({
            user: tenantScope.tenantId,
            projectId: id,
        });
        
        if (workLogCount > 0) {
            throw AppError.badRequest(
                'Impossibile eliminare progetto con worklog associati'
            );
        }
    }
}
\end{lstlisting}

\subsection{Esempio 3: Utilizzo con Paginazione}

\begin{lstlisting}
// Controller
const { data: goals, meta } = await goalService.paginate(
    req.tenantScope,
    { status: 'active' },
    { page: 1, limit: 20 }
);

// Risposta API
res.json({
    success: true,
    data: goals,
    pagination: meta, // { page, limit, total, totalPages, hasMore }
});
\end{lstlisting}

\section{Pattern e Principi}

\subsection{Template Method Pattern}

\textbf{Definizione}: Definisce lo scheletro di un algoritmo, delegando alcuni passi alle sottoclassi.

\textbf{Implementazione in BaseService}:
\begin{itemize}
    \item \textbf{Template (Metodo Base)}: \texttt{create()}, \texttt{update()}, \texttt{delete()}
    \item \textbf{Hook (Metodi Override)}: \texttt{beforeCreate()}, \texttt{afterUpdate()}, etc.
\end{itemize}

\textbf{Vantaggi}:
\begin{itemize}
    \item Codice comune in un solo posto
    \item Estensibilità senza modificare codice base
    \item Controllo del flusso centralizzato
\end{itemize}

\subsection{Repository Pattern (Semplificato)}

\textbf{Definizione}: Astrae l'accesso ai dati, nascondendo i dettagli di implementazione.

\textbf{Implementazione in BaseService}:
\begin{itemize}
    \item Controller chiama \texttt{goalService.find()}, non sa che usa MongoDB
    \item Se domani si cambia database, basta modificare BaseService
\end{itemize}

\subsection{Principi SOLID}

\begin{enumerate}
    \item \textbf{[S] Single Responsibility}: BaseService ha una sola responsabilità: CRUD con multi-tenancy
    \item \textbf{[O] Open/Closed}: Aperto per estensione (hook), chiuso per modifica
    \item \textbf{[L] Liskov Substitution}: Qualsiasi sottoclasse può sostituire BaseService
    \item \textbf{[I] Interface Segregation}: Metodi granulari, usa solo quelli necessari
    \item \textbf{[D] Dependency Inversion}: Dipende da astrazione (Model interface), non da implementazione
\end{enumerate}

\section{Conclusioni}

BaseService è un componente \textbf{critico} per la sicurezza e la manutenibilità del sistema. La sua implementazione di multi-tenancy esplicita previene vulnerabilità comuni (IDOR, mass assignment) e garantisce isolamento completo dei dati.

\textbf{Punti Chiave}:
\begin{itemize}
    \item ✅ Sicurezza by default con fail-fast
    \item ✅ DRY: riduce codice duplicato del 70-80\%
    \item ✅ Estensibilità tramite lifecycle hooks
    \item ✅ Consistenza: tutti i service si comportano allo stesso modo
    \item ⚠️ Possibili miglioramenti: validazione input, error handling standardizzato, escaping regex
\end{itemize}

\textbf{Raccomandazioni}:
\begin{enumerate}
    \item Implementare validazione input (ObjectId, limit, skip)
    \item Standardizzare error handling con AppError
    \item Escapare caratteri speciali in \texttt{search()}
    \item Considerare TypeScript per type safety
    \item Documentare requirement di indici MongoDB
\end{enumerate}

\vspace{1cm}

\textit{Documento generato il \today}

\end{document}
